%% Ce fichier est inclus à la fin de la section « Utiliser des modèles »
%% du chapitre 3 (03creerQuestionnaire.tex), juste après le modèle de champs.

\subsection{Créer un modèle de projet}
\label{modeleProjet}

Contrairement au modèle de champs, le modèle de projet englobe l'ensemble du projet : tous ses formulaires, toutes les pages de chaque formulaire et tout le code de contrôle (que nous verrons au chapitre~\ref{codeValidation}). C'est ce modèle que l'on envoie aux enquêteurs pour qu'ils disposent tous exactement du même questionnaire. Pour créer un modèle de projet à partir de notre projet COVID19 :

\begin{enumerate}
	\item Dans l'explorateur de projets, faire un clic droit sur le nom du projet, ici \var{COVID19}. Un menu contextuel apparaît ;
	\item Sélectionner \kw{Enregistrer le projet en tant que modèle} (\emph{Save Project as Template}). La fenêtre du même nom s'ouvre ;
	\item Taper le nom du modèle, par exemple \var{ModeleCOVID19}, dans le champ \kw{Nom du modèle} ;
	\item Cliquer sur \kw{OK}.
\end{enumerate}

\subsection{Envoyer un modèle de projet à un collègue}

Une fois le modèle créé, il apparaît dans l'explorateur de projets, dans la rubrique \kw{Modèles}, sous \var{Projects}. Il ne reste plus qu'à retrouver le fichier pour le joindre à un courriel ou le copier sur une clé USB :

\begin{enumerate}
	\item Repérer le modèle \var{ModeleCOVID19} sous \var{Projects} dans l'explorateur de projets. Il faudra peut-être faire défiler le volet vers le bas pour le trouver ;
	\item Faire un clic droit sur \var{ModeleCOVID19}. Un menu contextuel apparaît ;
	\item Sélectionner \kw{Ouvrir le dossier conteneur} (\emph{Open Containing Folder}). Une fenêtre de l'Explorateur de fichiers de Windows s'ouvre ;
	\item Repérer le fichier \var{ModeleCOVID19.xml} dans la liste. Ce fichier est le modèle : vous pouvez le joindre à un courriel avec votre messagerie habituelle ou le copier sur une clé USB.
\end{enumerate}

\subsection{Charger un modèle de projet reçu d'un collègue}

Mettons-nous maintenant à la place d'un enquêteur qui vient de recevoir le fichier \var{ModeleCOVID19.xml} par courriel. Pour créer le projet à partir de ce modèle :

\begin{enumerate}
	\item Enregistrer la pièce jointe sur le bureau de l'ordinateur ;
	\item Ouvrir le module \module{créer questionnaires} d'Epi Info 7 ;
	\item Aller dans \kw{Fichier > Obtenir un modèle\ldots} Une fenêtre Windows d'ouverture de fichier apparaît ;
	\item Aller sur le bureau et sélectionner le fichier reçu ;
	\item Cliquer sur \kw{Ouvrir}. La fenêtre \kw{Nouveau projet à partir d'un modèle} s'ouvre ;
	\item Sélectionner le modèle que vous venez de charger, s'il ne l'est pas déjà ;
	\item Si nécessaire, indiquer un nom de projet dans la zone \kw{Nom} ;
	\item Cliquer sur \kw{OK} pour créer le projet.
\end{enumerate}

L'enquêteur dispose alors d'un projet identique au vôtre, prêt pour la saisie des données.

\begin{mybox}[colbacktitle=TIP]{\faLightbulb{} CONSEIL}
Créez et envoyez le modèle de projet seulement après avoir testé le questionnaire et son code de contrôle (chapitres~\ref{codeValidation} et~\ref{saisirDonnees}). Tous les enquêteurs travailleront ainsi sur la même version, et leurs données pourront être fusionnées sans difficulté.
\end{mybox}
